% =============================================================
%  Appendix A2 -- Data Construction
%  Label: app:data
% =============================================================

\section{Data Construction}\label{app:data}

This appendix documents the construction of the analysis
dataset from raw sources, the coding decisions made at each
stage, and the sample attrition from the raw CFPS to the
estimation sample.

\subsection{Province-Level CPI Data}

The National Bureau of Statistics publishes monthly consumer
price sub-indices at the province level.
I collect the meat and poultry sub-index, the grain sub-index,
the eggs sub-index, and the headline CPI index for all 31
mainland provinces, municipalities, and autonomous regions
from January 2016 to December 2025.
The data are sourced from the NBS monthly statistical releases
and cross-checked against the CEIC China Premium Database.

Each observation is reported as a \emph{previous month $= 100$}
index: a value of 105 indicates that the price level in the
reference month is 5 percent higher than in the \emph{preceding}
month.
The NBS publishes CPI sub-indices in multiple bases
(previous month $= 100$, same month of preceding year $= 100$,
and fixed base); I use the month-on-month series throughout
because it permits exact construction of cumulative price
changes over arbitrary calendar windows without requiring a
year-on-year-to-monthly inversion.

I convert each monthly observation to a log change,
$\Delta\log(\mathrm{CPI}_{pm}) = \log(\mathrm{CPI}_{pm}/100)$.
Because the raw index uses the previous month as base,
$\mathrm{CPI}_{pm}/100$ equals the gross monthly price change,
and its logarithm is the exact month-on-month log growth rate.
For values close to 100 the approximation is accurate to
second order; for extreme values (such as meat CPI reaching
138 during the ASF peak) the log transformation compresses
the upper tail relative to a linear percentage change.

The treatment variable for the meat shock is the cumulative
sum of monthly log changes between the midpoints of
consecutive CFPS fieldwork periods:
\[
  \mathrm{MeatShock}_{p,\,t-1\to t}
  = \sum_{m \in \mathcal{M}(t-1,t)}
    \Delta\log(\mathrm{MeatCPI}_{pm}).
\]
The fieldwork midpoints are approximated as July of each wave
year (2016, 2018, 2020, 2022), based on the typical CFPS
survey timeline.
The grain shock and eggs shock are constructed identically
from their respective sub-indices.
All three commodity shocks vary at the province-wave level
and enter the regressions as right-hand-side variables.

\subsection{CFPS Panel Construction}

The China Family Panel Studies is a nationally representative
longitudinal survey conducted by the Institute of Social
Science Survey at Peking University.
I merge the household-level data across waves 2016, 2018,
2020, and 2022.
The key variables are the inflation expectation item
(qp201 in most waves), household demographic characteristics,
and the province of residence.

The inflation expectation variable is coded as
$\mu_{it} \in \{-1, 0, 1\}$, where $-1$ indicates the
household expects prices to fall, $0$ indicates prices will
stay the same, and $1$ indicates prices will rise.
The original CFPS coding uses numeric values (typically
1, 2, 3 or 1, 2, 5 depending on the wave), which I recode
to the $\{-1, 0, 1\}$ ordinal scale.
I verify the recoding against the CFPS codebook for each wave
to ensure consistency.

Province-of-residence codes present a harmonization challenge.
Across waves, the CFPS uses different variable names:
\texttt{provcd} in some waves, \texttt{provcd10} in the 2010
baseline, \texttt{provcd12} in the 2012 follow-up, and
variants in later waves.
All province codes follow the GB/T 2260 national standard but
differ in the number of trailing digits.
I standardise all codes to the two-digit level (e.g., Beijing
$= 11$, Shanghai $= 31$, Guangdong $= 44$) and merge on the
standardised code.
Provinces that changed administrative boundaries during the
sample period (none did at the provincial level between 2016
and 2022) would require additional harmonization, but this
issue does not arise.

The province code is used to match each household to
province-level CPI data.
The merge is many-to-one: all households within a province
at a given wave receive the same province-wave meat shock.
The CFPS covers 25 of the 31 mainland provinces in the public-use
data; the remaining provinces are either not sampled or
suppressed for confidentiality.
I use all available provinces, which yields coverage of
approximately 95 percent of the national population.

\subsection{Forecast Error Construction}

The forecast error is defined as
\[
  \mathrm{FE}_{ipt}
  = \pi^{\text{headline}}_{p,\,t\to t+1}
    - \tilde{\mu}_{ipt},
\]
where $\pi^{\text{headline}}_{p,\,t\to t+1}$ is the
annualised realised headline CPI inflation in province $p$
over the next inter-wave period, and $\tilde{\mu}_{ipt}$ is
the household's ordinal expectation scaled to CPI units.

The forward CPI is computed as follows.
For each province and wave, I take the cumulative sum of
monthly headline CPI log changes from the midpoint of wave
$t$ to the midpoint of wave $t+1$.
This cumulative change is then annualised by dividing by the
number of years between midpoints (approximately 2 for
biennial waves).
The annualisation ensures that the forecast error is
expressed in annual percentage-point units, comparable across
wave pairs of slightly different length.

For the scaling of the ordinal expectation, I map
$\{-1, 0, 1\}$ to $\{-2, 0, 2\}$ percentage points in the
baseline specification.
This mapping reflects the assumption that a household
reporting ``prices will rise'' expects approximately 2 percent
annual inflation, a household reporting ``stay the same''
expects zero, and a household reporting ``prices will fall''
expects approximately 2 percent deflation.
These values are chosen to be symmetric and to approximate
the unconditional mean of headline CPI inflation during the
sample period (which averages approximately 2 percent).
Appendix~\ref{app:scaling} reports sensitivity to alternative
scaling choices, including $\{-1, 0, 1\}$ (unit ordinal),
$\{-4, 0, 4\}$ (wide), and $\{-1, 0, 3\}$ (asymmetric).

The forecast-error sample requires that forward realised CPI
be available.
For the 2022 wave, the forward CPI requires data through
mid-2024, which is available in the current data extract.
For any wave whose forward CPI is not yet observed at the
time of data construction, the observation is dropped from
the forecast-error estimation but retained for the
expectation equation (Eq~1), which does not use forward CPI.

\subsection{Sample Attrition}\label{app:attrition}

The raw CFPS household panel contains approximately 130,000
household-wave observations across the four waves (2016, 2018,
2020, 2022).
The estimation sample is substantially smaller due to several
restrictions.

The first restriction is non-missing inflation expectations.
Approximately 15 to 20 percent of households do not respond
to the inflation expectation question (item qp201) in a given
wave, either because the question was not administered
(routing differences across waves) or because the respondent
refused or answered ``don't know.''
This restriction reduces the sample to approximately 105,000
observations.

The second restriction is a valid province code that matches
the CPI data.
Households in provinces not covered by the CFPS public-use
data, or with missing province codes, are dropped.
This restriction has a modest effect, reducing the sample to
approximately 100,000 observations.

The third restriction is non-missing meat-shock values.
Because the NBS meat CPI data begin in January 2016, the
earliest wave for which the meat shock can be constructed is
2016 (using the 2014--2016 backward link).
Households surveyed before 2016, if any remain after wave
selection, are dropped.
After this restriction, the full estimation sample for the
expectation equation (Eq~1) contains 90,133 household-wave
observations.

The fourth restriction applies only to the forecast-error
equation (Eq~3): forward realised headline CPI must be
available.
For wave $t$, this requires CPI data through the midpoint of
wave $t+1$.
At the time of data construction, forward CPI is available
for the 2016, 2018, and 2020 waves (using data through
mid-2022, mid-2024, and projected mid-2024 respectively).
The 2022 wave lacks a complete forward realisation at the
time of the most recent data pull, reducing the
forecast-error sample to 69,533 observations.

The fifth restriction applies to specifications with
demographic controls.
The age, education, and urban-status variables are drawn from
the CFPS household and individual questionnaires.
The urban indicator is missing for a substantial fraction of
observations in the 2018, 2020, and 2022 waves due to
changes in the questionnaire routing.
When all three demographic controls are required
simultaneously, the sample shrinks to approximately 22,121
observations.
The main results are reported with and without demographic
controls to show that the forecast-error finding is not
driven by the sample composition change.

\subsection{Persistence Estimation}

The AR(1) persistence of meat-price shocks is estimated from
the monthly province-level data.
The estimating equation is
\[
  \Delta\log(\mathrm{MeatCPI}_{pm})
  = \alpha_p
  + \rho\,\Delta\log(\mathrm{MeatCPI}_{p,m-1})
  + \varepsilon_{pm},
\]
estimated by OLS with province fixed effects and standard
errors clustered at the province level.
The sample covers January 2016 to December 2025 across 31
provinces, yielding 3,689 monthly observations after
dropping the initial month for each province.
The estimated coefficient is $\hat{\rho} = 0.480$
(s.e. = 0.010, $p < 0.001$).

The AR(1) specification is chosen for parsimony: it provides
a single persistence parameter that maps directly into the
signal-extraction model.
Higher-order specifications (AR(2) or AR(12) to capture
seasonal patterns) would allow richer dynamics but would
complicate the calibration without changing the qualitative
conclusion, because the rational bound depends on the total
persistence over the forecast horizon, not on the lag
structure.
The province fixed effects absorb persistent differences in
the level of meat-price inflation across provinces, ensuring
that $\hat{\rho}$ captures within-province temporal
persistence rather than cross-province level differences.

